% Teto: 500 palavras (orçamento de 28/10, tarefa 99 do painel).
% Texto: seção 1 do .docx do PR 60, migrado sem correção.
\section{Introdução}
\label{sec:introducao}

Redes de casas e de pequenos negócios, como um restaurante, reúnem na mesma rede televisores, câmeras, sensores e assistentes de voz, dispositivos da Internet das Coisas (IoT), e em geral não têm equipe de segurança. Nelas, um aparelho comprometido pode atacar os outros sem sair da rede. É o cenário do CICIoT2023, conjunto de dados em que 33 ataques, agrupados em sete categorias, são executados por dispositivos IoT maliciosos contra outros dispositivos IoT, numa rede de laboratório com 105 equipamentos \cite{neto2023}.

Os sistemas de detecção de intrusão na rede, ou NIDS (Network Intrusion Detection System), baseados em aprendizado de máquina apontam que o tráfego é ataque e de que categoria, mas param no alerta: entender o que aconteceu, escolher a resposta e aplicá-la continua manual e pede um especialista, justamente o que essas redes não têm. As plataformas comerciais de detecção e resposta de rede, ou NDR (Network Detection and Response), juntam a resposta ao alerta \cite{darktrace2024}, mas são produtos fechados, voltados a empresas. Este trabalho trata dessa lacuna com um sistema aberto de detecção e resposta a intrusões na rede, um NIDR (Network Intrusion Detection and Response).

Para a resposta, o trabalho usa agentes baseados em modelos de linguagem de grande escala, os LLM (Large Language Model), que alternam raciocínio e chamadas de ferramenta \cite{yao2023}. Eles leem o incidente, explicam o que o classificador viu e propõem uma ação. O Model Context Protocol (MCP) padroniza a ligação entre o modelo e as ferramentas \cite{hou2026}, e é por ele que os agentes chegam ao classificador e ao catálogo de ações. Levar dados de rede a um modelo tem riscos próprios: uma instrução escondida no tráfego pode mudar o comportamento do modelo \cite{greshake2023}, e um modelo pago na nuvem cobra por chamada e recebe dados da rede.

O objetivo geral é desenvolver e avaliar o SentryIoT, um sistema que detecta ataques em redes IoT com um classificador Random Forest e responde a eles com três agentes de LLM, ligados ao classificador por um servidor MCP e sujeitos a aprovação humana nas ações de risco alto.

Os objetivos específicos são: (i) extrair do tráfego bruto, em arquivo pcap, as 39 features do CICIoT2023, com um extrator calibrado contra o dos autores do dataset; (ii) treinar e avaliar o classificador em dados regerados dos pcaps, em janela única de 100 quadros, com controles em que o pcap inteiro fica fora do treino; (iii) construir um servidor MCP que aplica a política de risco e protege os endereços da rede sem depender do modelo; (iv) construir os agentes de triagem, decisão e execução numa arquitetura híbrida, em que só o agente de decisão roda na nuvem e nenhum endereço da rede sai do ambiente; e (v) comparar as configurações tudo local, tudo na nuvem e híbrida em qualidade da recomendação, custo em tokens e tempo até a remediação.

O público-alvo são as redes de casas e de pequenos negócios. O motivo principal é o tamanho da rede do dataset: o CICIoT2023 é uma rede de laboratório única, com 105 dispositivos e os atacantes do lado de dentro, o que corresponde a uma casa ou a um pequeno comércio, e não à rede de uma empresa. Nesse tamanho, uma caixa só captura o tráfego e aplica as respostas (seção~\ref{sec:captura}); quem não tem equipe de segurança é quem mais precisa de um agente que explique o incidente e proponha a resposta; e o custo de tokens pesa mais num orçamento pequeno (seção~\ref{sec:mcp}). A implantação numa empresa fica para os trabalhos futuros.

As contribuições do trabalho são: (1) o classificador treinado em dados regerados dos pcaps do CICIoT2023, em janela única e com o rótulo decidido pela presença do atacante na janela; (2) a medida de dois problemas do dataset, a dependência do modelo em relação ao tamanho da janela e as janelas rotuladas como ataque sem nenhum quadro do atacante; (3) o sistema de três agentes ligados por MCP, com as garantias de segurança no servidor e aprovação humana das ações de risco alto; e (4) a comparação entre as configurações local, nuvem e híbrida em qualidade, custo e tempo.

O artigo está organizado assim: a seção~\ref{sec:referencial} apresenta o referencial teórico e os trabalhos correlatos; a seção~\ref{sec:metodologia}, a metodologia, da captura do tráfego à avaliação dos agentes; a seção~\ref{sec:resultados}, os resultados e as limitações; e a seção~\ref{sec:conclusao}, as conclusões e os trabalhos futuros.
